\section{Setting and Measurement Protocol}
\label{sec:setting}

All experiments share the test speakers and the target labels of each iteration, so differences between methods are paired within an iteration.

\subsection{Corpora and Labels}
\label{sec:setting-data}

We use seven groups from six corpora (Table~\ref{tab:data}): the improvised and the scripted sessions of IEMOCAP~\cite{busso2008iemocap}, CMU-MOSEI~\cite{zadeh2018mosei}, MELD~\cite{poria2019meld}, AFEW~\cite{dhall2012afew}, RAVDESS~\cite{livingstone2018ravdess}, and CREMA-D~\cite{cao2014cremad}.
The two IEMOCAP session types form separate groups because all ten speakers read the same lines in the scripted sessions, which lets a text classifier recognize the script instead of the emotion.
Each utterance carries one of four labels: neutral, happy (merged with excited in IEMOCAP), sad, or angry; we exclude frustrated (IEMOCAP) and calm (RAVDESS).
A CMU-MOSEI utterance receives an emotion when exactly one of its six intensity ratings is at least $0.5$ and neutral when none is, so the neutral class of CMU-MOSEI is defined by the absence of rated emotions.

Audio exists in all seven groups, face video in six (we did not extract face features for RAVDESS), and transcripts in four (AFEW has none, and RAVDESS and CREMA-D read fixed sentences).
We call A (audio), V (face), and T (text) the single-modality configurations and AV, VT, AT, and AVT the fused configurations, which use only the utterances where all their modalities exist.
Speaker identity defines every split.
CMU-MOSEI has no speaker annotation; its channel identity is known for \NUM{data-mosei-chan-pct}\% of the videos, and the video identifier stands in for the speaker elsewhere, so its speaker separation is incomplete.

\begin{table}[t]
\caption{Groups, available modalities (A: audio, V: face, T: text), speaker groups, and utterances per emotion after filtering (counts for the audio features).}
\label{tab:data}
\centering
\footnotesize
\setlength{\tabcolsep}{3pt}
\begin{tabular}{@{}llrrrrrr@{}}
\toprule
Group & Mod. & Spk. & Utt. & Neu. & Hap. & Sad & Ang. \\
\midrule
IEMOCAP impro. & AVT & \NUM{data-spk-impro} & \NUM{data-n-impro-wavlm} & \NUM{data-impro-neu} & \NUM{data-impro-hap} & \NUM{data-impro-sad} & \NUM{data-impro-ang} \\
IEMOCAP script. & AVT & \NUM{data-spk-script} & \NUM{data-n-script-wavlm} & \NUM{data-script-neu} & \NUM{data-script-hap} & \NUM{data-script-sad} & \NUM{data-script-ang} \\
CMU-MOSEI & AVT & \NUM{data-spk-mosei}\textsuperscript{a} & \NUM{data-n-mosei-wavlm} & \NUM{data-mosei-neu} & \NUM{data-mosei-hap} & \NUM{data-mosei-sad} & \NUM{data-mosei-ang} \\
MELD & AVT & \NUM{data-spk-meld} & \NUM{data-n-meld-wavlm} & \NUM{data-meld-neu} & \NUM{data-meld-hap} & \NUM{data-meld-sad} & \NUM{data-meld-ang} \\
AFEW & AV & \NUM{data-spk-afew} & \NUM{data-n-afew-wavlm} & \NUM{data-afew-neu} & \NUM{data-afew-hap} & \NUM{data-afew-sad} & \NUM{data-afew-ang} \\
RAVDESS & A & \NUM{data-spk-ravdess} & \NUM{data-n-ravdess-wavlm} & \NUM{data-ravdess-neu} & \NUM{data-ravdess-hap} & \NUM{data-ravdess-sad} & \NUM{data-ravdess-ang} \\
CREMA-D & AV & \NUM{data-spk-cremad} & \NUM{data-n-cremad-wavlm} & \NUM{data-cremad-neu} & \NUM{data-cremad-hap} & \NUM{data-cremad-sad} & \NUM{data-cremad-ang} \\
\midrule
Total & & & \NUM{data-n-total-wavlm} & & & & \\
\bottomrule
\end{tabular}

\vspace{2pt}
\raggedright\textsuperscript{a}Channels where known, otherwise video identifiers.
\end{table}

\subsection{Frozen Features and the Linear Classifier}
\label{sec:setting-features}

No encoder is fine-tuned on the seven groups; the first encoder of each modality carries the analysis, and a second encoder tests whether the conclusions persist.
Audio features average the final layer of WavLM-Large~\cite{chen2022wavlm} (second: HuBERT-Large~\cite{hsu2021hubert}) over time, and text features average the final layer of RoBERTa-base~\cite{liu2019roberta} (second: DeBERTa-v3-base~\cite{he2023debertav3}) over tokens.
Face features average over the utterance the embeddings of the largest detected face in each sampled frame, computed by a ResNet-50 trained for facial expression recognition~\cite{ryumina2022facial} (second: the image encoder of CLIP ViT-L/14~\cite{radford2021clip}); the supplementary material gives the extraction details.
A fused configuration concatenates the single-modality features after standardizing each dimension.

The learner is multinomial logistic regression on standardized features with the penalty $\|W\|_F^2/(2C)$ on the weights and none on the intercepts, solved with L-BFGS~\cite{liu1989lbfgs} in double precision to a largest absolute gradient of $10^{-6}$.
Class weights inversely proportional to class frequency make it target the unweighted average recall (UAR), our main metric, whose chance level is 25\%.
A frozen encoder with a linear head keeps the decision rule open to analysis, and with a few labels per emotion fine-tuning can distort the pretrained features~\cite{kumar2022lpft}.

\subsection{Splits and Label Budget}
\label{sec:setting-splits}

For each target group, we split the speakers into five folds, and iteration $r$ tests on fold $r \bmod 5$, which holds \NUM{data-test-pct-mean}\% of the utterances on average (range \NUM{data-test-pct-lo}--\NUM{data-test-pct-hi}\%).
The remaining utterances form the pool, from which $K$ labeled utterances per emotion are drawn at random ($K=25$ means $100$ labels, and $K=0$ means none).
The draws are nested across $K$ and identical for every source, because the seed depends only on the target and the iteration.
The rest of the pool is available without labels, no method uses test utterances for training or selection, and the test speakers are also removed from the source, which matters only between the two IEMOCAP groups.

\subsection{Three Recipes and the Gain Map}
\label{sec:setting-map}

For every ordered pair of a target and a source group, each iteration trains three classifiers: \emph{target-only} on the $4K$ target labels, \emph{source-only} on all labeled source utterances, and \emph{pooled} on both.
The \emph{gain} of a source is the pooled UAR minus the target-only UAR in the same iteration.
The \emph{gain map} reports UAR for $K \in \{1, 5, 10, 25, 50\}$ and the full pool over \NUM{map-reps-a} iterations, for \NUM{map-pairs-a} ordered pairs in A, \NUM{map-pairs-v} in V and AV, and \NUM{map-pairs-t} in each configuration with text.

Three-fold speaker cross-validation selects $C \in \{0.01, 0.1, 1, 10\}$ for the target-only and the source-only classifiers ($C=1$ when a fold lacks an emotion), and the pooled classifier takes the $C$ of the source-only classifier of the same iteration, which uses no target label.
Validating $C$ on the target labels fails for small $K$, because with one label per emotion the validation scores tie in most iterations and the tie rule weakens the penalty; earlier cross-corpus work also fixes hyperparameters in this case~\cite{gideon2021addog}.
With the full pool, this rule overstates the loss of pooling for audio and changes it little for the other configurations; the supplementary material reports the map with cross-validated $C$.

\subsection{Threshold and Direction of a Pairwise Decision}
\label{sec:setting-model}

We analyze a four-class classifier through its six pairwise decisions.
For emotions $k$ and $l$, let $z = (\mathbf{w}_k - \mathbf{w}_l)^{\top}\mathbf{x} + (b_k - b_l)$ be the difference of their scores, and on the target utterances of $k$ and $l$ let $\mu_k$, $\mu_l$ be the class means of $z$ and $\sigma$ its pooled within-class standard deviation.
The \emph{separation} $d = (\mu_k - \mu_l)/\sigma$ depends only on the direction $\mathbf{w}_k - \mathbf{w}_l$, and the \emph{threshold offset} $o = -(\mu_k + \mu_l)/(2\sigma)$, the position of the boundary $z=0$ relative to the midpoint, can be set to zero by the intercepts alone.
If $z$ is Gaussian within each class, the pairwise accuracy (the mean of the two recalls) is
\begin{equation}
\label{eq:twocloud}
a(d, o) = \tfrac{1}{2}\left[\Phi\!\left(\tfrac{d}{2} + o\right) + \Phi\!\left(\tfrac{d}{2} - o\right)\right],
\end{equation}
where $\Phi$ is the standard normal distribution function.
Let $\Delta$ be the separation along the target's best direction, estimated by a shrinkage linear discriminant on one half of the target speakers and measured on the other half, averaged over the two halves; this reference serves only the analysis.
Against it, the loss of a classifier splits into a \emph{threshold loss} $a(d, 0) - a(d, o)$, which moving the threshold to the midpoint recovers, and a \emph{direction loss} $a(\Delta, 0) - a(d, 0)$, which remains; we average both over the six pairs, in points of pairwise accuracy.

\subsection{Hyperparameter Selection and Reporting}
\label{sec:setting-selection}

The hyperparameters of Part~II and of the competitors, which $K$ target labels cannot select, are selected with pseudo-targets, as in leave-one-domain-out selection~\cite{gulrajani2021domainbed}: for a target $T$, each source group $c$ from another corpus serves in turn as the target of the method built from the groups outside the corpora of $T$ and $c$, with $K$ labels per emotion drawn from $c$.
The selected value maximizes the UAR on the test speakers of $c$, averaged over $c$, the first five iterations, and the label budgets each method specifies.
No label of $T$ beyond its $K$ draws enters the selection; the \emph{oracle}, the best value on the test speakers of $T$, serves only to measure the loss due to selection.

We report UAR differences in percentage points, averaged over iterations within a pair (Part~I) or a target (Part~II and Section~\ref{sec:comparison}) and then over pairs or targets.
Intervals are 95\% $t$ intervals of the paired differences over iterations; they reflect the draw of labels and test speakers, not the resampling of speakers or corpora.
The ten IEMOCAP speakers yield only five (improvised) or seven (scripted) distinct test folds, so where we decide whether a source helps an IEMOCAP target, we first average the iterations that share a fold.
